\documentclass[12pt,a4paper]{article}
\usepackage[utf8]{inputenc}
\usepackage[T1]{fontenc}
\usepackage[polish]{babel}
\usepackage{amsmath}
\usepackage{amssymb}
\usepackage{geometry}
\usepackage{booktabs}
\usepackage{array}
\usepackage{longtable}
\usepackage{xcolor}
\usepackage{hyperref}
\usepackage{multirow}
\usepackage{colortbl}
\usepackage{float}

\geometry{margin=2cm}

\hypersetup{colorlinks=true, linkcolor=blue!70!black}

% ---- colours for topic rows ----
\definecolor{rowfoto}{RGB}{220,235,252}
\definecolor{rowzwierc}{RGB}{220,252,230}
\definecolor{rowsocz}{RGB}{255,245,210}
\definecolor{rowodbl}{RGB}{252,230,220}
\definecolor{rowrozs}{RGB}{240,220,252}
\definecolor{headercolor}{RGB}{50,80,140}

% ====================================================================
\title{\textbf{Fotonika -- Zestawienie oznaczeń i jednostek}\\[4pt]
       \large Ściągawka do kolokwium}
\author{}
\date{\today}
% ====================================================================

\begin{document}
\maketitle

\begin{center}
\textit{Plik zawiera wyłącznie zestawienie wszystkich symboli fizycznych
 używanych na kolokwium: co dany symbol oznacza i w jakiej jednostce się
 go podaje.  Szczegółowe wyjaśnienia i wyprowadzenia -- patrz \texttt{opracowanie.pdf}.}
\end{center}

\bigskip
\tableofcontents
\newpage

% ====================================================================
\section{Fotometria}
% ====================================================================

\begin{longtable}{>{\bfseries}p{1.6cm} p{6.5cm} p{3.2cm} p{3.5cm}}
\toprule
\rowcolor{headercolor}
\color{white}Symbol &
\color{white}Nazwa wielkości &
\color{white}Jednostka &
\color{white}Uwagi \\
\midrule
\endfirsthead
\toprule
\rowcolor{headercolor}
\color{white}Symbol &
\color{white}Nazwa wielkości &
\color{white}Jednostka &
\color{white}Uwagi \\
\midrule
\endhead
\bottomrule
\endfoot

\rowcolor{rowfoto}
$I$     & Światłość źródła światła
        & \textbf{kandela [cd]}
        & Jednostka podstawowa SI \\

\rowcolor{rowfoto}
$\Phi$  & Strumień świetlny
        & \textbf{lumen [lm]}
        & $1\,\mathrm{lm} = 1\,\mathrm{cd\cdot sr}$ \\

\rowcolor{rowfoto}
$E$     & Natężenie oświetlenia (iluminancja)
        & \textbf{luks [lx]}
        & $1\,\mathrm{lx} = 1\,\mathrm{lm/m^2}$ \\

\rowcolor{rowfoto}
$\Omega$ & Kąt bryłowy
        & \textbf{steradian [sr]}
        & $\Omega_{\text{pełny}} = 4\pi\,\mathrm{sr}$ \\

\rowcolor{rowfoto}
$A$     & Pole powierzchni
        & \textbf{metr kwadratowy [m$^2$]}
        & \\

\rowcolor{rowfoto}
$r$     & Odległość od źródła do powierzchni
        & \textbf{metr [m]}
        & \\

\rowcolor{rowfoto}
$\theta$ & Kąt między promieniem a normalną
          do powierzchni
        & \textbf{stopień [°] lub radian [rad]}
        & Używany w prawie Lamberta \\

\rowcolor{rowfoto}
$L$     & Luminancja (jasność)
        & \textbf{[cd/m$^2$]}
        & Świetlistość powierzchni \\

\rowcolor{rowfoto}
$K_m$   & Maks. skuteczność świetlna
        & \textbf{lumen na wat [lm/W]}
        & $K_m \approx 683\,\mathrm{lm/W}$ przy $\lambda=555\,\mathrm{nm}$ \\

\rowcolor{rowfoto}
$V(\lambda)$ & Krzywa czułości widmowej oka
        & bezwymiarowa ($0\leq V\leq 1$)
        & Maks. dla $\lambda=555\,\mathrm{nm}$ \\

\rowcolor{rowfoto}
$P$     & Moc promieniowania
        & \textbf{wat [W]}
        & \\

\rowcolor{rowfoto}
$\eta$  & Skuteczność świetlna
        & \textbf{[lm/W]}
        & $\Phi = \eta \cdot P$ \\

\end{longtable}

% ====================================================================
\section{Zwierciadła płaskie i sferyczne}
% ====================================================================

\begin{longtable}{>{\bfseries}p{1.6cm} p{6.5cm} p{3.2cm} p{3.5cm}}
\toprule
\rowcolor{headercolor}
\color{white}Symbol &
\color{white}Nazwa wielkości &
\color{white}Jednostka &
\color{white}Uwagi \\
\midrule
\endfirsthead
\toprule
\rowcolor{headercolor}
\color{white}Symbol &
\color{white}Nazwa wielkości &
\color{white}Jednostka &
\color{white}Uwagi \\
\midrule
\endhead
\bottomrule
\endfoot

\rowcolor{rowzwierc}
$f$     & Ogniskowa zwierciadła
        & \textbf{metr [m]}
        & $f>0$ -- wklęsłe; $f<0$ -- wypukłe \\

\rowcolor{rowzwierc}
$R$     & Promień krzywizny zwierciadła
        & \textbf{metr [m]}
        & $f = R/2$ \\

\rowcolor{rowzwierc}
$d_o$   & Odległość \textbf{przedmiotu} od zwierciadła
        & \textbf{metr [m]}
        & $d_o > 0$ -- przedmiot realny \\

\rowcolor{rowzwierc}
$d_i$   & Odległość \textbf{obrazu} od zwierciadła
        & \textbf{metr [m]}
        & $d_i>0$ -- obraz rzeczywisty;
          $d_i<0$ -- pozorny \\

\rowcolor{rowzwierc}
$m$     & Powiększenie liniowe
        & bezwymiarowe
        & $m = -d_i/d_o$;
          $m>0$ -- obraz prosty;
          $m<0$ -- odwrócony \\

\rowcolor{rowzwierc}
$h_o$   & Wysokość (rozmiar) przedmiotu
        & \textbf{metr [m]}
        & \\

\rowcolor{rowzwierc}
$h_i$   & Wysokość (rozmiar) obrazu
        & \textbf{metr [m]}
        & $m = h_i / h_o$ \\

\end{longtable}

% ====================================================================
\section{Soczewki}
% ====================================================================

\begin{longtable}{>{\bfseries}p{1.6cm} p{6.5cm} p{3.2cm} p{3.5cm}}
\toprule
\rowcolor{headercolor}
\color{white}Symbol &
\color{white}Nazwa wielkości &
\color{white}Jednostka &
\color{white}Uwagi \\
\midrule
\endfirsthead
\toprule
\rowcolor{headercolor}
\color{white}Symbol &
\color{white}Nazwa wielkości &
\color{white}Jednostka &
\color{white}Uwagi \\
\midrule
\endhead
\bottomrule
\endfoot

\rowcolor{rowsocz}
$f$     & Ogniskowa soczewki
        & \textbf{metr [m]}
        & $f>0$ skupiająca; $f<0$ rozpraszająca \\

\rowcolor{rowsocz}
$D$     & Zdolność skupiająca (moc optyczna)
        & \textbf{dioptria [D]} $= \mathrm{m}^{-1}$
        & $D = 1/f$ \\

\rowcolor{rowsocz}
$d_o$   & Odległość \textbf{przedmiotu} od soczewki
        & \textbf{metr [m]}
        & $d_o>0$ -- po lewej stronie soczewki \\

\rowcolor{rowsocz}
$d_i$   & Odległość \textbf{obrazu} od soczewki
        & \textbf{metr [m]}
        & $d_i>0$ -- rzeczywisty (prawa str.);
          $d_i<0$ -- pozorny (lewa str.) \\

\rowcolor{rowsocz}
$m$     & Powiększenie liniowe
        & bezwymiarowe
        & $m = -d_i/d_o$ \\

\rowcolor{rowsocz}
$n_L$   & Współczynnik załamania materiału soczewki
        & bezwymiarowy ($n\geq 1$)
        & \\

\rowcolor{rowsocz}
$n_0$   & Współczynnik załamania ośrodka otoczenia
        & bezwymiarowy
        & Powietrze: $n_0\approx 1{,}000$ \\

\rowcolor{rowsocz}
$R_1$   & Promień krzywizny \textbf{pierwszej} powierzchni
          soczewki
        & \textbf{metr [m]}
        & $R_1>0$: środek krzywizny po prawej \\

\rowcolor{rowsocz}
$R_2$   & Promień krzywizny \textbf{drugiej} powierzchni
          soczewki
        & \textbf{metr [m]}
        & $R_2<0$: środek krzywizny po lewej
          (typowe dla bikonweksowej) \\

\rowcolor{rowsocz}
$h_o$   & Rozmiar przedmiotu
        & \textbf{metr [m]}
        & \\

\rowcolor{rowsocz}
$h_i$   & Rozmiar obrazu
        & \textbf{metr [m]}
        & \\

\end{longtable}

% ====================================================================
\section{Odbicie i załamanie światła}
% ====================================================================

\begin{longtable}{>{\bfseries}p{1.6cm} p{6.5cm} p{3.2cm} p{3.5cm}}
\toprule
\rowcolor{headercolor}
\color{white}Symbol &
\color{white}Nazwa wielkości &
\color{white}Jednostka &
\color{white}Uwagi \\
\midrule
\endfirsthead
\toprule
\rowcolor{headercolor}
\color{white}Symbol &
\color{white}Nazwa wielkości &
\color{white}Jednostka &
\color{white}Uwagi \\
\midrule
\endhead
\bottomrule
\endfoot

\rowcolor{rowodbl}
$n$     & Bezwzględny współczynnik załamania ośrodka
        & bezwymiarowy ($n\geq 1$)
        & $n = c/v$ \\

\rowcolor{rowodbl}
$n_1$   & Wsp. załamania ośrodka \textbf{padania}
        & bezwymiarowy
        & \\

\rowcolor{rowodbl}
$n_2$   & Wsp. załamania ośrodka \textbf{załamania}
        & bezwymiarowy
        & \\

\rowcolor{rowodbl}
$c$     & Prędkość światła w próżni
        & \textbf{metr na sekundę [m/s]}
        & $c \approx 3\times10^8\,\mathrm{m/s}$ \\

\rowcolor{rowodbl}
$v$     & Prędkość fazowa fali w ośrodku
        & \textbf{metr na sekundę [m/s]}
        & $v \leq c$ \\

\rowcolor{rowodbl}
$\theta_1$ & Kąt \textbf{padania} (od normalnej)
        & \textbf{stopień [°] lub radian [rad]}
        & Mierzony od normalnej do powierzchni \\

\rowcolor{rowodbl}
$\theta_r$ & Kąt \textbf{odbicia} (od normalnej)
        & \textbf{stopień [°] lub radian [rad]}
        & Prawo odbicia: $\theta_r = \theta_1$ \\

\rowcolor{rowodbl}
$\theta_2$ & Kąt \textbf{załamania} (od normalnej)
        & \textbf{stopień [°] lub radian [rad]}
        & \\

\rowcolor{rowodbl}
$\theta_c$ & Kąt graniczny (krytyczny) całkowitego
             wewnętrznego odbicia
        & \textbf{stopień [°] lub radian [rad]}
        & $\sin\theta_c = n_2/n_1$;
          wymaga $n_1 > n_2$ \\

\end{longtable}

% ====================================================================
\section{Rozszczepienie światła}
% ====================================================================

\begin{longtable}{>{\bfseries}p{1.6cm} p{6.5cm} p{3.2cm} p{3.5cm}}
\toprule
\rowcolor{headercolor}
\color{white}Symbol &
\color{white}Nazwa wielkości &
\color{white}Jednostka &
\color{white}Uwagi \\
\midrule
\endfirsthead
\toprule
\rowcolor{headercolor}
\color{white}Symbol &
\color{white}Nazwa wielkości &
\color{white}Jednostka &
\color{white}Uwagi \\
\midrule
\endhead
\bottomrule
\endfoot

\rowcolor{rowrozs}
$\lambda$ & Długość fali światła
        & \textbf{nanometr [nm]} lub \textbf{metr [m]}
        & $1\,\mathrm{nm} = 10^{-9}\,\mathrm{m}$;
          zakres VIS: $380$--$780\,\mathrm{nm}$ \\

\rowcolor{rowrozs}
$\nu$   & Częstotliwość fali światła
        & \textbf{teraherc [THz]} lub \textbf{herc [Hz]}
        & $1\,\mathrm{THz} = 10^{12}\,\mathrm{Hz}$;
          $\nu = c/\lambda$ \\

\rowcolor{rowrozs}
$c$     & Prędkość światła w próżni
        & \textbf{metr na sekundę [m/s]}
        & $c \approx 3\times10^8\,\mathrm{m/s}$ \\

\rowcolor{rowrozs}
$n(\lambda)$ & Współczynnik załamania zależny od
               długości fali (dyspersja)
        & bezwymiarowy
        & Rośnie gdy $\lambda$ maleje
          (dyspersja normalna) \\

\rowcolor{rowrozs}
$V_d$   & Liczba Abbego (współczynnik dyspersji)
        & bezwymiarowy
        & Duże $V_d$ = mała dyspersja \\

\rowcolor{rowrozs}
$n_d$   & Wsp. załamania dla linii $d$ He
          ($\lambda_d = 587{,}6\,\mathrm{nm}$, żółta)
        & bezwymiarowy
        & Linia referencyjna \\

\rowcolor{rowrozs}
$n_F$   & Wsp. załamania dla linii $F$ H
          ($\lambda_F = 486{,}1\,\mathrm{nm}$, niebieska)
        & bezwymiarowy
        & \\

\rowcolor{rowrozs}
$n_C$   & Wsp. załamania dla linii $C$ H
          ($\lambda_C = 656{,}3\,\mathrm{nm}$, czerwona)
        & bezwymiarowy
        & \\

\rowcolor{rowrozs}
$\varphi$ & Kąt łamiący (wierzchołkowy) pryzmatu
        & \textbf{stopień [°]}
        & Kąt między obiema powierzchniami
          łamiącymi pryzmatu \\

\rowcolor{rowrozs}
$\delta$ & Kąt odchylenia promienia przez pryzmat
        & \textbf{stopień [°]}
        & $\delta = (\theta_1 + \theta_2') - \varphi$ \\

\rowcolor{rowrozs}
$\delta_{\min}$ & Minimalny kąt odchylenia promienia
                  w pryzmacie
        & \textbf{stopień [°]}
        & Symetryczny bieg promienia;
          służy do pomiaru $n_p$ \\

\rowcolor{rowrozs}
$n_p$   & Współczynnik załamania materiału pryzmatu
        & bezwymiarowy
        & $n_p = \dfrac{\sin\!\left(\frac{\varphi+\delta_{\min}}{2}\right)}%
                       {\sin\!\left(\frac{\varphi}{2}\right)}$ \\

\rowcolor{rowrozs}
$n_0$   & Współczynnik załamania ośrodka wokół
          pryzmatu (zwykle powietrze)
        & bezwymiarowy
        & $n_0 \approx 1$ \\

\rowcolor{rowrozs}
$\theta_1$ & Kąt padania na pierwszą powierzchnię
             pryzmatu
        & \textbf{stopień [°]}
        & Mierzony od normalnej \\

\rowcolor{rowrozs}
$\theta_1'$ & Kąt załamania na pierwszej powierzchni
              (wewnątrz pryzmatu)
        & \textbf{stopień [°]}
        & $n_0\sin\theta_1 = n_p\sin\theta_1'$ \\

\rowcolor{rowrozs}
$\theta_2$ & Kąt padania na drugą powierzchnię
             (wewnątrz pryzmatu)
        & \textbf{stopień [°]}
        & $\theta_1' + \theta_2 = \varphi$ \\

\rowcolor{rowrozs}
$\theta_2'$ & Kąt wyjścia z drugiej powierzchni
              pryzmatu
        & \textbf{stopień [°]}
        & $n_p\sin\theta_2 = n_0\sin\theta_2'$ \\

\end{longtable}

% ====================================================================
\section{Zestawienie zbiorcze -- wszystkie symbole}
% ====================================================================

\begin{longtable}{>{\bfseries}p{2.0cm} p{7.0cm} p{3.2cm} p{2.5cm}}
\toprule
\rowcolor{headercolor}
\color{white}Symbol &
\color{white}Opis &
\color{white}Jednostka &
\color{white}Temat \\
\midrule
\endfirsthead
\toprule
\rowcolor{headercolor}
\color{white}Symbol &
\color{white}Opis &
\color{white}Jednostka &
\color{white}Temat \\
\midrule
\endhead
\bottomrule
\endfoot

\rowcolor{rowfoto}
$I$            & Światłość        & \textbf{cd}                & Fotometria \\
\rowcolor{rowfoto}
$\Phi$         & Strumień świetlny & \textbf{lm}               & Fotometria \\
\rowcolor{rowfoto}
$E$            & Natężenie oświetlenia & \textbf{lx}            & Fotometria \\
\rowcolor{rowfoto}
$\Omega$       & Kąt bryłowy      & \textbf{sr}                & Fotometria \\
\rowcolor{rowfoto}
$L$            & Luminancja       & \textbf{cd/m}$^2$          & Fotometria \\
\rowcolor{rowfoto}
$K_m$          & Maks. skuteczność świetlna & \textbf{lm/W}    & Fotometria \\
\rowcolor{rowzwierc}
$f$            & Ogniskowa zwierciadła / soczewki & \textbf{m}  & Zwierciadła, Soczewki \\
\rowcolor{rowzwierc}
$R$            & Promień krzywizny & \textbf{m}                & Zwierciadła, Soczewki \\
\rowcolor{rowzwierc}
$d_o$          & Odległość przedmiotu & \textbf{m}             & Zwierciadła, Soczewki \\
\rowcolor{rowzwierc}
$d_i$          & Odległość obrazu  & \textbf{m}               & Zwierciadła, Soczewki \\
\rowcolor{rowzwierc}
$m$            & Powiększenie liniowe & bezwymiarowe            & Zwierciadła, Soczewki \\
\rowcolor{rowzwierc}
$h_o, h_i$     & Rozmiar przed./obrazu & \textbf{m}           & Zwierciadła, Soczewki \\
\rowcolor{rowsocz}
$D$            & Zdolność skupiająca & \textbf{D} (dioptria)   & Soczewki \\
\rowcolor{rowsocz}
$n_L$          & Wsp. zał. soczewki & bezwymiarowy             & Soczewki \\
\rowcolor{rowsocz}
$R_1, R_2$     & Promienie krzywizny powierzchni & \textbf{m}  & Soczewki \\
\rowcolor{rowodbl}
$n$, $n_1$, $n_2$ & Współczynnik załamania & bezwymiarowy      & Odbicie/załamanie \\
\rowcolor{rowodbl}
$c$            & Prędkość światła w próżni & \textbf{m/s}       & Odbicie/załamanie, Dyspersja \\
\rowcolor{rowodbl}
$v$            & Prędkość fali w ośrodku & \textbf{m/s}         & Odbicie/załamanie \\
\rowcolor{rowodbl}
$\theta_1$     & Kąt padania      & \textbf{° lub rad}          & Odbicie/załamanie \\
\rowcolor{rowodbl}
$\theta_r$     & Kąt odbicia      & \textbf{° lub rad}          & Odbicie/załamanie \\
\rowcolor{rowodbl}
$\theta_2$     & Kąt załamania    & \textbf{° lub rad}          & Odbicie/załamanie \\
\rowcolor{rowodbl}
$\theta_c$     & Kąt graniczny c.w.o. & \textbf{° lub rad}      & Odbicie/załamanie \\
\rowcolor{rowrozs}
$\lambda$      & Długość fali     & \textbf{nm lub m}           & Dyspersja, Pryzmat \\
\rowcolor{rowrozs}
$\nu$          & Częstotliwość    & \textbf{THz lub Hz}         & Dyspersja, Pryzmat \\
\rowcolor{rowrozs}
$V_d$          & Liczba Abbego    & bezwymiarowy                & Dyspersja \\
\rowcolor{rowrozs}
$n_d, n_F, n_C$ & Wsp. zał. dla linii spektralnych & bezwymiarowe & Dyspersja \\
\rowcolor{rowrozs}
$\varphi$      & Kąt łamiący pryzmatu & \textbf{°}             & Pryzmat \\
\rowcolor{rowrozs}
$\delta$       & Kąt odchylenia w pryzmacie & \textbf{°}         & Pryzmat \\
\rowcolor{rowrozs}
$\delta_{\min}$ & Minimalny kąt odchylenia & \textbf{°}          & Pryzmat \\
\rowcolor{rowrozs}
$n_p$          & Wsp. zał. pryzmatu & bezwymiarowy              & Pryzmat \\

\end{longtable}

\bigskip
\begin{center}
  \textbf{Legenda kolorów:}
  \quad
  \colorbox{rowfoto}{\strut\ Fotometria\ }
  \quad
  \colorbox{rowzwierc}{\strut\ Zwierciadła\ }
  \quad
  \colorbox{rowsocz}{\strut\ Soczewki\ }
  \quad
  \colorbox{rowodbl}{\strut\ Odbicie/Załamanie\ }
  \quad
  \colorbox{rowrozs}{\strut\ Rozszczepienie\ }
\end{center}

\end{document}
